\documentclass[aspectratio=169,10pt]{beamer}
\usetheme{Madrid}
\usecolortheme{seahorse}
\usefonttheme{professionalfonts}
\usepackage[utf8]{inputenc}
\usepackage[T1]{fontenc}
\usepackage[spanish]{babel}
\usepackage{booktabs}
\usepackage{amsmath}
\usepackage{listings}
\usepackage{xcolor}
\definecolor{codebg}{RGB}{246,248,250}
\definecolor{deepblue}{RGB}{0,68,120}
\setbeamercolor{structure}{fg=deepblue}
\setbeamercolor{frametitle}{fg=white,bg=deepblue}
\setbeamercolor{title}{fg=deepblue}
\setbeamertemplate{navigation symbols}{}
\setbeamertemplate{footline}[frame number]
\lstset{basicstyle=\ttfamily\small, backgroundcolor=\color{codebg}, frame=single,
  columns=fullflexible, showstringspaces=false}

\title{Punto fijo para filtros FIR}
\subtitle{De enteros binarios a formatos Q y cálculos en Arduino}
\author{Francisco Carlos Calderon Ph.D}
\date{\today}

\begin{document}

\begin{frame}
  \titlepage
\end{frame}

\begin{frame}{Propósito de esta sesión}
  Al terminar, se debe poder:
  \begin{itemize}
    \item interpretar un entero como un número con punto binario fijo;
    \item elegir un formato con signo, $W$ bits y $F$ bits fraccionarios;
    \item calcular el resultado y el tamaño de una suma o multiplicación;
    \item explicar cada escala usada por un FIR cuantizado.
  \end{itemize}
  \vfill
  \alert{Idea central:} el procesador guarda enteros. Nosotros decidimos qué
  factor de escala representa cada entero.
\end{frame}

\begin{frame}{Recordatorio: un entero binario}
  Con cuatro bits sin signo se representan enteros entre $0$ y $15$:
  \[
    \texttt{1101}_2 = 1\cdot 2^3+1\cdot 2^2+0\cdot 2^1+1\cdot 2^0=13.
  \]
  Los pesos de los bits son:
  \[
    \begin{array}{c|cccc}
    \text{posición} & 3 & 2 & 1 & 0\\ \hline
    \text{peso} & 8 & 4 & 2 & 1
    \end{array}
  \]
  Un entero no lleva punto decimal. El punto fijo asigna algunos bits a pesos
  menores que uno.
\end{frame}

\begin{frame}{El primer ejemplo de punto fijo}
  Guardemos el mismo entero $13$, pero acordemos dividirlo siempre entre $4$:
  \[
    x=\frac{13}{4}=3.25.
  \]
  En binario se puede leer como:
  \[
    \texttt{1101}_2 \quad\longrightarrow\quad \texttt{11.01}_2=3.25.
  \]
  El hardware sigue almacenando \texttt{1101}. El punto no ocupa memoria: es
  una regla de interpretación compartida por el programa.
\end{frame}

\begin{frame}{Regla general}
  Si un entero almacenado se llama $r$ y se reservan $F$ bits fraccionarios:
  \[
    \boxed{x=\frac{r}{2^F}=r\,2^{-F}}
  \]
  \begin{exampleblock}{Ejemplo}
    $r=45$ y $F=4$:\quad $x=45/16=2.8125$.
  \end{exampleblock}
  El menor incremento representable se llama \alert{LSB}:
  \[
    \mathrm{LSB}=2^{-F}.
  \]
  Más bits fraccionarios dan pasos más pequeños, pero reducen el rango si $W$
  permanece constante.
\end{frame}

\begin{frame}{W, F, QF y QI}
  Usaremos una notación explícita para evitar ambigüedades:
  \begin{center}
  \begin{tabular}{ll}
    $W$ & número total de bits almacenados\\
    $F$ o QF & bits fraccionarios\\
    QI & bits enteros sin contar el signo\\
    signo & un bit adicional cuando el formato es con signo
  \end{tabular}
  \end{center}
  Para un formato con signo:
  \[
    W=1+QI+QF.
  \]
  Ejemplo: con signo, $W=8$, $F=4$. Entonces $QI=3$ y escribiremos
  \texttt{Q3.4}.
\end{frame}

\begin{frame}{Formato con signo}
  Los formatos con signo usan complemento a dos. Para $W=8$ y $F=4$:
  \[
    \texttt{0010.1101}_2=45/16=2.8125
  \]
  \[
    \texttt{1110.1100}_2=(-20)/16=-1.25
  \]
  El bit más a la izquierda tiene peso negativo. En este caso sus pesos son:
  \[
    -2^3,\;2^2,\;2^1,\;2^0,\;2^{-1},\;2^{-2},\;2^{-3},\;2^{-4}.
  \]
\end{frame}



\begin{frame}{Rango de un formato}
  Para un formato \alert{sin signo} de $W$ bits y $F$ fraccionarios:
  \[
    0 \leq x \leq \frac{2^W-1}{2^F}.
  \]
  Para un formato \alert{con signo}:
  \[
    -\frac{2^{W-1}}{2^F}\leq x\leq
    \frac{2^{W-1}-1}{2^F}.
  \]
  \begin{center}
  \begin{tabular}{lccc}
  \toprule
  formato & mínimo & máximo & LSB\\ \midrule
  sin signo, $W=8,F=4$ & $0$ & $15.9375$ & $1/16$\\
  con signo, $W=8,F=4$ & $-8$ & $7.9375$ & $1/16$\\
  \bottomrule
  \end{tabular}
  \end{center}
\end{frame}

\begin{frame}{La notación Q puede cambiar entre libros}
  Algunos autores cuentan el bit de signo dentro de la parte entera y otros no.
  Por eso \texttt{Q1.15} puede significar cosas distintas según la fuente.
  \vfill
  Se recomienda declarar siempre:
  \[
    \boxed{\texttt{signed},\quad W,\quad F}
  \]
  y mostrar QI y QF sólo como ayuda pedagógica, así se encuentra en varias aplicaciones en línea.
  \begin{exampleblock}{En MATLAB}
    Un formato con signo, $W=16$ y $F=15$ se escribe
    \texttt{numerictype(1,16,15)}.
  \end{exampleblock}
\end{frame}

\begin{frame}{Codificar y decodificar un número}
  Ejemplo: formato con signo $W=8$, $F=4$.
  \begin{enumerate}
    \item Para guardar $x=2.75$, multiplicar por $2^F$:
      \[r=2.75\cdot 16=44.\]
    \item Almacenar el entero $r=44$:
      \[44=\texttt{00101100}_2.\]
    \item Para recuperar el valor físico:
      \[x=44/16=2.75.\]
  \end{enumerate}
  Si $x\,2^F$ no es entero, hay que escoger un redondeo.
\end{frame}

\begin{frame}{Cuantización y redondeo}
  En formato $W=8$, $F=4$, el valor $0.30$ no se representa exactamente:
  \[
    0.30\cdot 16=4.8.
  \]
  \begin{center}
  \begin{tabular}{lccc}
  \toprule
  método & entero almacenado & valor representado & error\\ \midrule
  hacia abajo & $4$ & $4/16=0.25$ & $-0.05$\\
  al más cercano & $5$ & $5/16=0.3125$ & $+0.0125$\\
  hacia cero & $4$ & $0.25$ & $-0.05$\\
  \bottomrule
  \end{tabular}
  \end{center}
  La elección debe formar parte de la especificación del filtro.
\end{frame}

\begin{frame}{Desbordamiento: saturar}
  Con signo, $W=4$, $F=0$, el rango que cabe en el registro es de $-8$ a $7$.
  \vspace{2mm}
  \textbf{Suma matemática:} $7+3=10$. El resultado correcto es mayor que el
  máximo representable ($+7$), por lo que no cabe en cuatro bits con signo.
  \begin{center}
  \begin{tabular}{lcc}
  \toprule
  operación & resultado matemático & valor guardado con saturación\\ \midrule
  $7+3$ & $10$ & $7$ (\texttt{0111})\\
  \bottomrule
  \end{tabular}
  \end{center}
  \begin{block}{Regla de saturación}
    Si el cálculo supera $+7$, se guarda $+7$. Si fuera menor que $-8$, se
    guardaría $-8$. El resultado queda limitado al rango del formato.
  \end{block}
  Aquí la suma real vale $10$, pero el registro conserva el máximo positivo:
  \texttt{0111}, que representa $+7$.
\end{frame}

\begin{frame}{Desbordamiento: wrap}
  Con el mismo formato con signo ($W=4$, $F=0$), el hardware conserva sólo
  cuatro bits del resultado de la suma:
  \[
    \begin{array}{rcl}
      \texttt{0111} & = & +7\\[-1mm]
      \texttt{+0011} & = & +3\\ \hline
      \texttt{1\,1010} & = & 10\quad\text{(resultado de cinco bits)}\\[-1mm]
      \texttt{1010} & & \text{(se descarta el bit que no cabe)}
    \end{array}
  \]
  \textbf{Interpretación de los cuatro bits guardados:} en complemento a dos,
  \texttt{1010} tiene un bit con peso $-8$ y otro con peso $+2$:
  \[
    \texttt{1010}_2=-8+2=-6.
  \]
  Wrap equivale a calcular módulo $2^4=16$: $10\bmod16=10$. La palabra binaria
  guardada sigue siendo \texttt{1010}; al leerla como número con signo se obtiene
  $-6$. Por eso wrap puede cambiar el signo del resultado.
\end{frame}

\begin{frame}{Tipos usuales en Arduino Uno}
  El Uno R3 usa el ATmega328P. En el compilador AVR habitual:
  \begin{center}
  \begin{tabular}{lcc}
  \toprule
  tipo & bits & rango con signo\\ \midrule
  \texttt{char} / \texttt{int8\_t} & 8 & $-128$ a $127$\\
  \texttt{short} / \texttt{int16\_t} & 16 & $-32768$ a $32767$\\
  \texttt{int} & 16 & $-32768$ a $32767$\\
  \texttt{long} / \texttt{int32\_t} & 32 & aproximadamente $\pm2.1\times10^9$\\
  \texttt{long long} / \texttt{int64\_t} & 64 & aproximadamente $\pm9.2\times10^{18}$\\
  \texttt{float} & 32 & punto flotante IEEE 754 de precisión simple\\
  \bottomrule
  \end{tabular}
  \end{center}
  El Uno tiene 2 kB de SRAM. Cada bit y cada arreglo importan.
\end{frame}

\begin{frame}[fragile]{No suponer tamaños de C/C++}
  El estándar permite que \texttt{int} cambie entre plataformas. Para exigir un
  ancho concreto, usar \texttt{<stdint.h>}:
  \begin{lstlisting}[language=C++]
#include <stdint.h>
#include <limits.h>

static_assert(sizeof(int16_t) == 2,
              "Se requieren 16 bits");
static_assert(CHAR_BIT == 8,
              "Se requieren bytes de 8 bits");

int16_t muestra;
int32_t acumulador;
  \end{lstlisting}
  \texttt{int16\_t} existe sólo si la plataforma ofrece exactamente 16 bits.
  \texttt{sizeof} devuelve bytes, no bits.
\end{frame}

\begin{frame}[fragile]{Multiplicación de enteros}
  Si $a$ usa $W_a$ bits y $b$ usa $W_b$ bits, el producto completo requiere:
  \[
    \boxed{W_p=W_a+W_b}
  \]
  bits como regla segura, tanto para formatos con signo como sin signo.
  \begin{exampleblock}{Ejemplo sin signo}
    $255\times 1023=260865$. Un operando de 8 bits y otro de 10 bits pueden
    producir hasta 18 bits de resultado.
  \end{exampleblock}
  Antes de multiplicar en C++, promover los operandos mediante casting ():
  \begin{lstlisting}[language=C++]
int32_t producto = (int32_t)a * (int32_t)b;
  \end{lstlisting}
\end{frame}

\begin{frame}{Multiplicación de números de punto fijo}
  Las escalas también se multiplican:
  \[
    \frac{r_x}{2^{F_x}}\,\frac{r_h}{2^{F_h}}
    =\frac{r_xr_h}{2^{F_x+F_h}}.
  \]
  Por tanto:
  \[
    \boxed{F_p=F_x+F_h}\qquad \boxed{W_p=W_x+W_h}.
  \]
  Ejemplo: entrada con signo $(W_x,F_x)=(16,15)$ y coeficiente con signo
  $(W_h,F_h)=(16,15)$:
  \[
    \text{producto: }(W_p,F_p)=(32,30).
  \]
\end{frame}

\begin{frame}{Suma de enteros y acumulación}
  Para sumar dos números de $W$ bits con la misma escala, se necesita un bit
  adicional en el peor caso:
  \[
    W_{\text{suma}}=W+1.
  \]
  Para sumar $N$ productos del mismo formato:
  \[
    \boxed{W_{\text{acc}}=W_p+\lceil\log_2(N)\rceil},\qquad F_{\text{acc}}=F_p.
  \]
  Si las escalas son distintas, primero hay que alinear los puntos binarios.
  Esa conversión puede introducir redondeo y overflow.
\end{frame}

\begin{frame}{Cadena numérica de un FIR}
  Para un FIR de 32 taps con entrada y coeficientes $(W,F)=(16,15)$:
  \[
  \begin{array}{rcl}
    x[n] &:& (16,15)\\
    h[k] &:& (16,15)\\
    x[n-k]h[k] &:& (32,30)\\
    \sum_{k=0}^{31}x[n-k]h[k] &:& (37,30)\\
    y[n] &:& (16,15)
  \end{array}
  \]
  Para pasar del acumulador $(37,30)$ a la salida $(16,15)$ se eliminan 15
  bits fraccionarios. El redondeo decide cómo eliminarlos y el modo de
  overflow decide qué ocurre si el resultado no cabe en 16 bits.
\end{frame}

\begin{frame}{Dos significados de normalizar}
  \textbf{1. Reescalar un número:} cambiar de $F=30$ a $F=15$ para guardar una
  salida. Es una operación de formato y requiere redondeo.
  \vfill
  \textbf{2. Normalizar la ganancia del filtro:} escoger los coeficientes para
  que la ganancia deseada se cumpla. En un pasobajo:
  \[
    H(0)=\sum_k h[k].
  \]
  Si el diseño ideal suma $1$, la ganancia DC ideal es uno. Después de
  cuantizar, hay que revisar:
  \[
    \widehat{H}(0)=\sum_k \widehat{h}[k].
  \]
\end{frame}

\begin{frame}{Ejemplo: la cuantización puede cambiar la ganancia}
  Coeficientes ideales de un pasobajo:
  \[
    h=[0.2,\;0.6,\;0.2],\qquad \sum h[k]=1.
  \]
  Con $F_h=3$, el paso es $1/8$. Al redondear:
  \[
    \widehat{h}=[2/8,\;5/8,\;2/8]=[0.25,\;0.625,\;0.25].
  \]
  \[
    \sum\widehat{h}[k]=1.125.
  \]
  La escala numérica sigue siendo correcta, pero el filtro ahora tiene ganancia
  DC de $1.125$. La herramienta debe mostrar ambas ganancias.
\end{frame}

\begin{frame}{Ejercicio manual}
  Formato de entrada: con signo, $W_x=8$, $F_x=4$.\\[4pt]
  Formato del coeficiente: con signo, $W_h=8$, $F_h=4$.\\[8pt]
  \begin{enumerate}
    \item Codifique $x=1.5$ y $h=-0.75$ como enteros raw.
    \item Multiplique los enteros raw.
    \item Indique $W_p$ y $F_p$.
    \item Decodifique el producto y compárelo con $1.5\times(-0.75)$.
  \end{enumerate}
  \vfill
  Respuesta esperada: $r_x=24$, $r_h=-12$, $r_p=-288$, $F_p=8$ y
  $-288/256=-1.125$.
\end{frame}

\begin{frame}{Ideas para recordar}
  \begin{itemize}
    \item Punto fijo guarda un entero y usa una escala $2^{-F}$.
    \item Declarar \texttt{signed}, $W$ y $F$ evita ambigüedades.
    \item Producto: sumar longitudes de palabra y bits fraccionarios.
    \item Acumulador FIR: añadir bits de guarda para las sumas.
    \item Cada conversión de formato necesita una política de redondeo y
    overflow.
    \item La ganancia del filtro se revisa después de cuantizar coeficientes.
  \end{itemize}
\end{frame}

\begin{frame}{Fuentes y material para continuar}
  \footnotesize
  \begin{itemize}
    \item Arduino UNO R3: ATmega328P, 16 MHz, 2 kB SRAM y 32 kB Flash.\\
    \url{https://docs.arduino.cc/hardware/uno-rev3}
    \item Arduino, tipos \texttt{float}: 32 bits y precisión simple.\\
    \url{https://docs.arduino.cc/language-reference/en/variables/data-types/float/}
    \item C/C++, \texttt{<stdint.h>} y los tipos de ancho fijo.\\
    \url{https://en.cppreference.com/c/types/integer}
    \item Laboratorio local: \texttt{fir\_fixed\_lab}, diseño, simulación y
    exportación compatible con Arduino.
  \end{itemize}
\end{frame}

\end{document}
